\section{The principal sector with both polar conditions retained}

Here the raw space is \(L^2(\mathbb R_+,du)\), identified with
the even full-line space with its corresponding normalization.
The even Fourier transform is
\[
 (\mathcal F\psi)(\xi)=2\int_0^\infty
                       \psi(u)\cos(2\pi u\xi)\,du .
\]
Write \(I\psi=\int_0^\infty\psi(u)\,du\).
The raw prolate projection \(\Pi_\Lambda\) retains exactly the
original even indices \(n\leq4\Lambda^2\); its raw time and
frequency radii are both \(\Lambda\).
All assertions below concern sufficiently large \(\Lambda\),
so the retained concentration values have the uniform lower
bound \(\nu\geq1/4\), and the first omitted value is at most
\(3/4\), as in AC-T29 and AC-T33.
Put \(c=\Lambda^2\), \(L=\log\Lambda\), and define the actual
finite constrained raw space by
\[
 \mathcal R_\Lambda^0
   =\operatorname{ran}\Pi_\Lambda
                 \cap\ker I\cap\ker\operatorname{ev}_0 .
 \tag{TP.1}
\]
Point evaluation in (TP.1) means evaluation of the smooth
interior prolate representative; the sharp endpoint causes
no ambiguity at zero.  This is a definition of the precise
archimedean finite family used in this section.
It is not an identification of two different adelic cutoff
constructions.

\begin{lemma}[AC-T42: corrected theta and the actual two-constraint projection]
\label{AC-T42}
For even Schwartz functions define
\[
 E_\zeta\psi(x)=\sqrt{x}\sum_{n\geq1}\psi(nx)
                         -x^{-1/2}I\psi .
 \tag{TP.2}
\]
Its closed Mellin realization is multiplication by
\(\zeta(1/2+i\tau)\) after the raw unitary
\((U\psi)(t)=e^{t/2}\psi(e^t)\), and it satisfies
\[
 E_\zeta\mathcal F=J E_\zeta,\qquad
 (JF)(x)=F(1/x),\qquad
 \|E_\zeta\psi\|_2\leq C\|\psi\|_{H_D^{5/16}} .
 \tag{TP.3}
\]
If a raw vector has time support in \([-\Lambda,\Lambda]\),
then its upper target tail is
\(-I\psi\,x^{-1/2}\) for \(x>\Lambda\).
If it has Fourier support in that interval, its lower
target tail is
\(-\psi(0)\sqrt{x}/2\) for \(x<\Lambda^{-1}\).

For every fixed smooth compactly supported even \(\psi\),
\(\Pi_\Lambda\psi\to\psi\) in raw and \(E_\zeta\)-graph
norm, and both \(I(\Pi_\Lambda\psi)\to I\psi\) and
\((\Pi_\Lambda\psi)(0)\to\psi(0)\).
Every fixed smooth compactly supported even vector satisfying
\(I\psi=\psi(0)=0\) is approximated, in raw norm and in the
\(E_\zeta\)-graph norm, by vectors of \(\mathcal R_\Lambda^0\).
There is also the following quantitative endpoint version.
For \(m\geq12\), let \(X_m\) be the full-line spatial \(H^m\)
space of even functions supported in \([-1,1]\).
For an \(X_m\)-valued Hilbert--Schmidt column operator \(\Psi\)
whose columns satisfy both constraints, the unitary raw
dilations \(D_\Lambda\Psi\), or their Fourier transforms,
can be approximated by actual columns
\(\Psi_\Lambda^0\in\mathcal R_\Lambda^0\) with
\[
 \|E_\zeta(\Psi_\Lambda^0-D_\Lambda\Psi)\|_2
 \leq C_m c^{-m/6+5/16}\|\Psi\|_{2,X_m}.
 \tag{TP.4}
\]
The identical inequality holds with
\(D_\Lambda\Psi\) replaced by
\(\mathcal F D_\Lambda\Psi\).
Here \(D_r\psi(u)=r^{-1/2}\psi(u/r)\).
For a fixed smooth profile either error decreases faster
than every power of \(c\).
\end{lemma}
\begin{proof}
Poisson summation and
\(\widehat\psi(0)=2I\psi\),
\(I\widehat\psi=\psi(0)/2\), give
\[
 \sqrt{x}\sum_{n\geq1}\psi(nx)-x^{-1/2}I\psi
 =x^{-1/2}\sum_{n\geq1}\widehat\psi(n/x)
                                      -\sqrt{x}\,\psi(0)/2.
 \tag{TP.5}
\]
This proves the Fourier identity and both asserted tails.
The Mellin identity follows first where the regularized
integral converges after subtraction of the displayed
integral term, and then by analytic continuation to
\(\operatorname{Re}s=1/2\).  More explicitly, splitting
the integral at \(x=1\), subtracting the leading Poisson
term on \((0,1)\), and integrating the subtracted power
separately gives
\(\zeta(s)\int_0^\infty\psi(u)u^{s-1}du\).
The terms in (TP.5) decay sufficiently rapidly at their
respective ends to justify this continuation on the
critical line.  The ordinary convexity bound for \(\zeta\)
there, with the harmless enlarged exponent \(5/16\), gives
the last estimate in (TP.3) by Plancherel.
These identities extend to the indicated closed domain.
For a band-supported vector the value at zero is continuous
in its band \(L^2\) norm; applying (TP.5) on smooth band
cores therefore also proves the stated lower-tail identity
for such vectors.  The upper-tail identity extends in the
same way on a fixed compact raw time interval, where \(I\)
is continuous.

Choose one even smooth \(\beta\), supported in
\(\{1/4<|u|<1/2\}\), with \(I\beta=1\), and set
\[
 b_{\Lambda,+}=D_\Lambda\beta,\qquad
 b_{\Lambda,-}=\tfrac12\mathcal F D_\Lambda\beta,\qquad
 \ell_{\Lambda,+}=I/\sqrt{\Lambda},\quad
 \ell_{\Lambda,-}=\operatorname{ev}_0/\sqrt{\Lambda}.
 \tag{TP.6}
\]
The matrix \((\ell_{\Lambda,i}(b_{\Lambda,j}))_{i,j}\)
is exactly the identity.  The theta norms of the two
columns are independent of \(\Lambda\), by (TP.3) and
dilation covariance.  Their selected projections
\(\widetilde b_{\Lambda,j}=\Pi_\Lambda b_{\Lambda,j}\)
satisfy, for every \(A>0\),
\[
 \|E_\zeta(\widetilde b_{\Lambda,j}-b_{\Lambda,j})\|_2
     =O_A(c^{-A}),\qquad
 A_\Lambda:=(\ell_{\Lambda,i}
                   (\widetilde b_{\Lambda,j}))_{i,j}
                         =1+O_A(c^{-A}).
 \tag{TP.7}
\]
Here is a verification of the constraints in (TP.7),
in addition to the graph estimate of AC-T38 or AC-S29.
The raw projection errors have every polynomial decay
rate.  On the time interval, Cauchy--Schwarz bounds the
normalized mean error by that raw norm.  The omitted
time tail of the Fourier profile has the same estimate
by its Schwartz decay.  On the interior time interval
the retained prolate derivative bound is
\(\|(\Pi_\Lambda f)'\|_2\leq C\Lambda\|f\|_2\).
The one-dimensional inequality
\[
 |v(0)|^2\leq C\|v\|_2\|v'\|_2+
                          C\Lambda^{-1}\|v\|_2^2
 \tag{TP.8}
\]
on \((-\Lambda,\Lambda)\) follows by averaging the
fundamental theorem of calculus for \(|v|^2\).
Combining an arbitrarily small raw error with this
coarse derivative bound proves the normalized
point-evaluation assertion in (TP.7).
Thus \(A_\Lambda^{-1}\) is uniformly bounded.

For any \(f\in\operatorname{ran}\Pi_\Lambda\), the exact
correction
\[
 {\cal C}_\Lambda f
   =f-(\widetilde b_{\Lambda,+},\widetilde b_{\Lambda,-})
                 A_\Lambda^{-1}
             \binom{\ell_{\Lambda,+}(f)}
                   {\ell_{\Lambda,-}(f)}
 \tag{TP.9}
\]
belongs to \(\mathcal R_\Lambda^0\).
This is a finite-dimensional linear map; no invariance
of the two constraints under \(\Pi_\Lambda\) is assumed.

For a dilated time profile from \(X_m\), the full-cutoff
raw approximation estimate is
\(O(c^{-m}\|\psi\|_{X_m})\).
In coordinates scaled to the unit time interval, the
derivative of its retained projection is
\(O(c\|\psi\|_{X_m})\).
Equation (TP.8) therefore bounds its normalized
point-evaluation error by
\(C c^{-m/2}\|\psi\|_{X_m}\); the normalized mean
error is \(O(c^{-m}\|\psi\|_{X_m})\).
For the Fourier profile its omitted raw time integral
is bounded by Cauchy--Schwarz against
\((1+|u|)^{-m}\), after the exact scaling
\(u\mapsto\Lambda u\); it is again \(O(c^{-m})\)
after mean normalization.  The needed raw projection
estimate does not assert that \(\Pi_\Lambda\) commutes
with Fourier transform.  Indeed, for a time-supported
vector \(f=\sum_j a_j\varphi_j\) in the complete
orthonormal prolate basis,
\[
 \begin{split}
 \|(1-\Pi_\Lambda)\mathcal Ff\|_2^2
   &=\sum_{j\ {\rm retained}}(1-\nu_j)|a_j|^2
                        +\sum_{j\ {\rm omitted}}|a_j|^2\\
   &\leq4\sum_j(1-\nu_j)|a_j|^2
                    =4\|(1-P)\mathcal Ff\|_2^2 .
 \end{split}
\]
The first equality uses
\(P\mathcal F P\varphi_j=i^j\sqrt{\nu_j}\varphi_j\),
up to the harmless Fourier-phase convention; the
inequality uses the first omitted value at most
\(3/4\).
Fourier conjugation reverses the dilation generator
and preserves every \(H_D^s\) norm, so the same
graph interpolation applies to this raw error.
Its derivative has the
same \(O(\Lambda)\) band bound, so (TP.8) gives the
same \(O(c^{-m/2})\) normalized value error.
These estimates also hold after summing squares
over any orthonormal input basis.
Apply (TP.9) to the retained projection.
The graph bound of AC-S29, whose proof uses only the
critical-line multiplier bound in (TP.3), is
\(C_m c^{-m/6+5/16}\|\Psi\|_{2,X_m}\).
The correction adds at most
\(C_m c^{-m/2}\|\Psi\|_{2,X_m}\), which is smaller.
This proves (TP.4) and its Fourier version.

For a fixed, undilated compact smooth vector,
the raw time tail is eventually zero and its Fourier
tail has every polynomial rate in \(\Lambda\).
The same graph interpolation and (TP.8) apply directly.
Its unnormalized mean error is at most
\(\sqrt{\Lambda}\) times its raw error, and its
unnormalized value error is bounded using the
coarse \(O(\Lambda)\) derivative bound in (TP.8).
The available arbitrary polynomial rates absorb
both factors.  This proves the unconstrained
fixed-profile assertion.  If both constraints vanish,
(TP.9) proves the constrained assertion as well.
The identical argument
applies to the smooth bulk \(P_\Lambda^pQ_\Lambda^p
P_\Lambda^pT\) of AC-T38 after its two constraints
have been corrected using the unprojected pair
in (TP.6): its remaining graph and constraint
errors are \(O_A(c^{-A}\|T\|_2)\), with any fixed
number of output derivatives included.
\end{proof}

\begin{lemma}[AC-T43: quantitative principal half-line approximation]
\label{AC-T43}
Assume \(\mathsf A\).
If a Hilbert--Schmidt column operator \(F\) into
\(L^2((-\infty,0])\) satisfies
\[
 \|F\|_2+\|F\|_{2,H^{1/4}}+\|tF\|_2\leq B,
 \tag{TP.10}
\]
with \(B\) sufficiently large, it can be approximated
within \(1/100\), in Hilbert--Schmidt norm, by
\(E_\zeta\Psi\), where every column of \(\Psi\) is
smooth, even, supported in a common compact subset
of \(\{0<|u|<1\}\), and satisfies \(I\psi=\psi(0)=0\).
For each fixed \(m\),
\[
                 \|\Psi\|_{2,X_m}\leq \exp(C_mB^{102}).
 \tag{TP.11}
\]
These columns, dilated to the upper endpoint or
Fourier-transformed to the lower endpoint, can be
replaced by actual columns in \(\mathcal R_\Lambda^0\)
with total target error at most \(1/50\), as soon as
\(\log\Lambda\geq C B^{102}\).
Changing these fixed tolerances changes only constants.
\end{lemma}
\begin{proof}
We give the modifications of AC-T41-halfline explicitly,
including the polar condition.  Set
\(A(s)=(s-1)\zeta(s)\).
Under \(\mathsf A\), it is analytic and zero-free on
\(\operatorname{Re}s>1/2\).
The Harnack argument of (TT.3), now applied to \(A\),
and the elementary rational factors below give the
same inverse bound with exponent \(C/\delta\).
At the disk center \(2+it\), its extra factor has
logarithm \(O(\log(2+|t|))\); this does not change
that argument.  In the far right half-plane the
Euler product and \(|(s-1)/(s+\delta-1)|\leq C\)
give the needed uniform bound directly.

Perform exactly the negative translation mollification
and far-left smoothing used in (TT.10)--(TT.11):
\(M=C_0B\), \(\eta=(C_0B)^{-4}\),
\(K=B^{10}\), \(\delta=B^{-50}\).
They produce \(F_0\) supported in
\([-2M-3\eta,-2\eta]\), within \(1/1000\) of \(F\),
with \(\|F_0\|_{2,H^1}\leq CB^5\) and
\[
 \|F_0\|_{2,H^n}
  \leq\exp\{C(n+1)(\log(n+2)+\log(B+2))\}.
 \tag{TP.12}
\]
All these estimates follow from (TP.10), the
translation bound in \(H^{1/4}\), and the first
log-moment bound, independently of the \(L\)-function.
For \(s=1/2+iz\), define
\[
 \widehat h_\delta(z)
       =\frac{(s-1)\widehat F_0(z-i\delta)}
                          {A(s+\delta)}.
 \tag{TP.13}
\]
This quotient is analytic for \(\operatorname{Im}z<\delta\).
The apparent principal pole has been removed by \(A\).
The same inverse bound and (TP.12), with
\(n\) a sufficiently large fixed multiple of
\(\delta^{-1}\), prove all the fixed Sobolev estimates
\(\exp(C_jB^{51})\) on the real line and on
\(\operatorname{Im}z=\delta/2\), exactly as in
(TT.13)--(TT.15).

For completeness the two contour uses have different
purposes.  Moving the line to \(z=x-iV\), \(V\to\infty\),
the numerator and its integrable \(x\)-derivatives are
bounded by a polynomial in \(V\) times \(e^{-2\eta V}\).
The far-right Euler bound controls the rationally
modified reciprocal in (TP.13).
The inverse transform contributes \(e^{-Vt}\).
For \(t>-\eta\) this tends to zero; vertical sides
vanish by sufficiently many of (TP.12)'s derivatives.
Thus \(h_\delta\) is supported in \(t\leq-\eta\).
Moving instead to \(\operatorname{Im}z=\delta/2\)
gives the negative-tail estimate
\[
 \|e^{-\delta t/2}h_\delta\|_{2,H^j}
                          \leq\exp(C_jB^{51}).
 \tag{TP.14}
\]

The theta multiplier of (TP.13) is
\[
 \widetilde M_\delta(t)
 =\frac{\zeta(s)(s-1)}{A(s+\delta)}
 =\frac{\xi(s)}{\xi(s+\delta)}
       \frac{s+\delta}{s}\,
       \pi^{-\delta/2}
       \frac{\Gamma((s+\delta)/2)}{\Gamma(s/2)},
 \qquad s=\tfrac12+it .
 \tag{TP.15}
\]
The function \(t\mapsto\xi(1/2+it)\) is real entire
of order one and, under \(\mathsf A\), has only real
zeros.  For the normalization and functional equation
see DLMF (25.4.3)--(25.4.4),
\url{https://dlmf.nist.gov/25.4}.
Its requisite entire growth bound also follows
directly from the same Poisson summation:
\[
 \xi(s)=\frac12+\frac{s(s-1)}2
       \int_1^\infty
       \left(\sum_{n\geq1}e^{-\pi n^2u}\right)
       \bigl(u^{s/2}+u^{(1-s)/2}\bigr)\,\frac{du}{u}.
\]
The integral converges locally uniformly in \(s\)
and is bounded on \(|s|\leq r\) by a gamma integral
with logarithm \(O(r\log(2+r))\).
Thus the order is at most one, which suffices for
the genus-one Hadamard factorization used here.
Conjugation and \(\xi(s)=\xi(1-s)\) give reality on
the critical line.  The Hadamard-product argument (TT.4)--(TT.6)
therefore applies with its full multiplicities.
The additional rational factor in (TP.15) is bounded
and differs from one by \(O(\delta)\) on the real
line; the gamma factor has the same bounds as before.
Consequently
\[
 |\widetilde M_\delta(t)|
       \leq C(2+|t|)^{\delta/2},\qquad
 \|\widetilde M_\delta-1\|_{L^2([-K,K])}
             \leq C K\log(2+K)\sqrt{\delta}.
 \tag{TP.16}
\]
There is no condition on the spacing of zeros.
The low-frequency row bound
\(\sup_t|\widehat F_0(t)|\leq C\sqrt M B\),
(TP.16), and the \(H^1\) bound in (TP.12) give,
respectively, the same errors
\(O(B^{-13.5}\log B)\) and \(O(B^{-3.75})\)
as in AC-T41-halfline.  The damping translation
costs \(O(\delta MB)\).
Increasing the lower threshold for \(B\) thus makes
\(\|E_\zeta U^{-1}h_\delta-F\|_2<1/200\).

There is now an exact additional fact:
\[
 \widehat h_\delta(-i/2)=0,\qquad
       I(U^{-1}h_\delta)
           =\int e^{t/2}h_\delta(t)\,dt=0 .
 \tag{TP.17}
\]
Indeed \(s=1\) at that argument, and the numerator
in (TP.13) vanishes while \(A(1+\delta)\ne0\).
The integral converges because of the proved negative
support and its \(L^2\) bound.
Cut \(h_\delta\) smoothly below \(-H\), where
\(H=CB^{102}\), using a transition of fixed width.
Equation (TP.14) makes its graph error arbitrarily
small and bounds the resulting spatial \(H^m\)
norm by \(\exp(C_mB^{102})\).
The mean row introduced by this cutoff has norm at
most \(e^{-H/2}\) times the discarded \(L^2\) norm.
Subtract that row times the fixed \(\beta\) in
(TP.6).  This restores \(I=0\) exactly, with
negligible graph and spatial cost.
Both functions vanish near raw zero, so their
value at zero is exactly zero as well.
This proves (TP.11) and the claimed approximation.
Finally apply (TP.4) with \(m=12\), whose exponent
is \(-27/16\), and then the exact constraint
correction (TP.9).  The exponential cost in (TP.11)
is absorbed when \(\log\Lambda\geq CB^{102}\).
The Fourier construction preserves both constraints
by the two identities preceding (TP.5), and (TP.3)
places its output at the opposite target endpoint.
\end{proof}

\begin{theorem}[AC-T44: the full tested trace for the actual finite principal \(S_0\) family]
\label{AC-T44}
Assume \(\mathsf A\).  Let \(Q_{\Lambda,0}\) be the
orthogonal projection in \(L^2(\mathbb R_+,dx/x)\) onto
\(E_\zeta\mathcal R_\Lambda^0\), with (TP.1) and the
original full even prolate cutoff.
Let \(S_\Lambda=1_{[\Lambda^{-1},\Lambda]}\) and
\[
 D_{\Lambda,0}
       =S_\Lambda-S_\Lambda Q_{\Lambda,0}S_\Lambda .
 \tag{TP.18}
\]
For every original norm test
\(h\in C_c^\infty(\mathbb R_+)\),
\[
 \operatorname{Tr}(Q_{\Lambda,0}V(h))
       =2\log\Lambda\,h(1)
          -\sum_{\zeta(1/2+i\gamma)=0}
                    m_\gamma\widehat h_+(\gamma)+o(1).
 \tag{TP.19}
\]
The corresponding positive defect measures converge
against every Schwartz Mellin-frequency test.
This theorem concerns the two-constraint subspace.
It does not include the two polar characters in the full
extension \(B_\Lambda/B_{\Lambda,0}\); that extension
requires its own explicit metric and trace calculation.
\end{theorem}
\begin{proof}
We give all the changes needed to the already proved
nonprincipal finite-image argument.  In particular neither
the principal polar tails nor the failure of
\(\Pi_\Lambda\) to preserve the two constraints is omitted.

\emph{1. The actual Gram matrix and tested tails.}
Take any raw orthonormal basis of \(\mathcal R_\Lambda^0\),
let \(E_0\) be its theta column operator, and put
\(G_\Lambda^{\rm pr}=E_0^*E_0\).
The multiplier \(\zeta(1/2+i\tau)\) is nonzero almost
everywhere, so this matrix is positive definite.
For every \(\kappa>0\),
\[
        \|(G_\Lambda^{\rm pr})^{-1}\|
                              \leq C_\kappa c^\kappa .
 \tag{TP.20}
\]
Here is the principal zero-free input to AC-T28.
The analytic zero-free function
\[
              F(s)=\frac{s-1}{s+1}\zeta(s),
                       \qquad \operatorname{Re}s>1/2,
 \tag{TP.21}
\]
has an analytic logarithm under \(\mathsf A\).
Its logarithm is uniformly bounded on the disks of
radius \(1/2\) centered at \(2+i\tau\), by the
Euler product and the displayed rational factor.
Its real part has the same \(O(\log(2+|\tau|))\)
upper bound on the larger disks used in (TB.32).
The Borel--Caratheodory and three-circles proof there
therefore gives, for fixed \(\delta,\epsilon>0\),
\[
        |\zeta(1/2+\delta+i\tau)^{-1}|
                          \leq C_{\delta,\epsilon}
                                       (1+|\tau|)^\epsilon.
 \tag{TP.22}
\]
Indeed \(1/\zeta=((s-1)/(s+1))/F\), and that rational
factor is bounded on the stated half-plane.
Each vector of \(\mathcal R_\Lambda^0\) has mean zero,
so its theta image vanishes above \(\Lambda\).
Thus both weighted functions in (TB.35)--(TB.37)
have the same upper support bound as before.
Their Mellin identity uses \(\zeta\) in place of \(L_\chi\).
The prolate derivative and fractional sharp-endpoint
estimates are unchanged for a subspace of the selected
raw range.  Applying (TB.33)--(TB.40) with (TP.22)
proves (TP.20), with exactly the arbitrarily small
exponent allowed there.  No bound for a different
truncated Gram matrix is being substituted here.

Let \(P=1_{[-\Lambda,\Lambda]}\) and
\(Q=\mathcal F^{-1}1_{[-\Lambda,\Lambda]}\mathcal F\)
be the sharp raw projections, and put
\[
         T_g=V(g)1_{(0,\Lambda^{-1})}E_\zeta .
 \tag{TP.23}
\]
For \(g\in\mathscr G=C_c^\infty(\mathbb R_+)+\mathbb C k\),
\(k(e^t)=\tfrac12e^{-|t|}\), this has a bounded
extension to the raw space.
The proof is the multiplier--half-line commutator
argument of AC-T25: \(\zeta\) and its first derivative
on the critical line have the same polynomial bounds
used there, while the multiplier of \(g\), or
\((1+\tau^2)^{-1}\) for \(k\), supplies the stated
decay.  This argument involves only the critical
line, which contains no pole of \(\zeta\).

The outside-time kernel, with input measure \(du\)
and output measure \(dy/y\), is now exactly
\[
 K_{\Lambda,g}(y,u)=u^{-1/2}
 \left\{\sum_{n>\Lambda u}n^{-1/2}g(ny/u)
       -\int_{\Lambda u}^{\infty}v^{-1/2}g(vy/u)\,dv
 \right\},\qquad u>\Lambda .
 \tag{TP.24}
\]
The integral is the mean subtraction in (TP.2), after
the same change of variables as in (TB.44).
Euler summation, with the bounded discrepancy
\(\lfloor v\rfloor-v\), bounds (TP.24) by
\[
       |K_{\Lambda,g}(y,u)|
          \leq C\frac{\sqrt y}{u}A_g(\Lambda y),
 \tag{TP.25}
\]
where \(A_g\) is precisely (TB.45).
One integration by parts proves this: the endpoint
and the integral of the absolute derivative of
\(v^{-1/2}g(vy/u)\) are at most
\((y/u)^{1/2}A_g(\Lambda y)\).
In particular
\[
                  \|T_g(1-P)\|_2^2\leq C_g\Lambda^{-2}.
 \tag{TP.26}
\]
The verification for \(k\) is unchanged because
\(A_k(s)=1\) for \(s\leq1\) and \(s^{-3/2}\) for
\(s\geq1\).

It would be incorrect to assert \(T_gQ=0\) on the
whole principal band space: its lower polar tail
depends on the value at zero.
Instead, for \(\psi\in\mathcal R_\Lambda^0\) put
\[
 v=(PQP|_{\operatorname{ran}\Pi_\Lambda})^{-1}\psi .
\]
This inverse has norm at most \(\nu^{-1}\).
The prolate equation gives \(PQv=\psi\), including
equality of their values at zero.
Since \(\psi(0)=0\), (TP.5) gives \(T_gQv=0\).
Consequently \(T_g\psi=-T_g(1-P)Qv\).
Summing over the raw orthonormal basis and using
(TP.26) proves
\[
 \|V(g)(1-S_\Lambda)E_0\|_2^2
                         \leq C_g\nu^{-2}\Lambda^{-2}.
 \tag{TP.27}
\]
The upper target tail is zero because the raw means
vanish.  Normalizing by the actual Gram matrix gives
\[
 \|V(g)(1-S_\Lambda)E_0(G_\Lambda^{\rm pr})^{-1/2}\|_2^2
                  \leq C_{g,\kappa}\Lambda^{-2+2\kappa}.
 \tag{TP.28}
\]
The Hilbert--Schmidt expansion (TB.54), the interval
energy bound, and the explicit \(k\)-factorization
and polarization in AC-T30 now prove
\[
 \operatorname{Tr}\bigl((Q_{\Lambda,0}
          -S_\Lambda Q_{\Lambda,0}S_\Lambda)V(h)\bigr)
                                      \longrightarrow0
 \tag{TP.29}
\]
for every original \(h\).  These are purely
operator identities and require no character sum
after (TP.28).

\emph{2. Every zero jet still gives the full lower bound.}
Fix \(s_0=1/2+i\gamma\), a zero of multiplicity \(m\).
For \(0\leq j<m\), the lower-tail functional has the
following bounded raw kernel:
\[
 \begin{split}
 H_{\gamma,j,\Lambda}(u)
  =\left.(\partial_s+\log\Lambda)^j
   \left[
    u^{s-1}\left(\zeta(s)-\sum_{n\leq\Lambda u}n^{-s}\right)
          +\frac{\Lambda^{1-s}}{1-s}
   \right]\right|_{s=s_0}.
 \end{split}
 \tag{TP.30}
\]
For a Schwartz raw vector, cut the defining theta
integral below at \(\varepsilon\) and interchange
the finite sum.  Euler summation gives
\(\sum_{n\leq u/\varepsilon}n^{-s}
 =\zeta(s)+(u/\varepsilon)^{1-s}/(1-s)
       +O_s((u/\varepsilon)^{-\operatorname{Re}s})\).
The power of \(\varepsilon\) cancels exactly against
the mean subtraction in (TP.2).  The remaining
constant is \(\Lambda^{1-s}/(1-s)\), with the sign
shown in (TP.30).  Differentiation yields the
logarithmic powers in the tail integral.
For small \(u<\varepsilon\) the omitted contribution
is \(O(\sqrt{\varepsilon}(1+|\log\varepsilon|)^{j+1})\);
for \(u\geq\varepsilon\), the Euler remainder gives
the same bound, with at most one further logarithm,
after integration against a bounded compact raw
vector.  This proves the identity without an
absolute infinite-sum interchange.

If \(u<\Lambda^{-1}\), all derivatives of the
\(\zeta\)-term in (TP.30) vanish, and the kernel is
bounded by \(C_{\gamma,j}\sqrt{\Lambda}\).
If \(u\geq\Lambda^{-1}\), Euler's bounded discrepancy,
differentiated in the combination
\(\partial_s+\log(\Lambda u)\), gives
\[
       |H_{\gamma,j,\Lambda}(u)|
                       \leq C_{\gamma,j}\Lambda^{-1/2}u^{-1}.
 \tag{TP.31}
\]
For example the Euler remainder is an endpoint
term plus an integral of a bounded discrepancy
against \(\partial_v v^{-s}\).
After differentiation its total variation is
\(O_{\gamma,j}((\Lambda u)^{-1/2})\), since the
inserted powers are \(\log(v/(\Lambda u))\);
the corresponding integrals converge.
It follows that the functional is bounded on raw
\(L^2\), and its restriction to \(u>\Lambda\)
has squared norm \(O_{\gamma,j}(\Lambda^{-2})\).
For a band vector with value zero at the origin,
its lower target tail is zero by (TP.5), so this
functional annihilates that vector.  Smooth band
cores and the bounded kernel extend this assertion
to the whole band space with that constraint.
The inverse-prolate argument following (TP.26)
therefore proves the analogue of (TP.28) for each
of these zero-tail functionals.

Each raw vector in (TP.1) has bounded variation on
its compact time interval.  Euler summation makes
its corrected theta image \(O(\sqrt{x})\) at zero,
and its mean-zero condition makes it vanish above
\(\Lambda\).  Hence its full differentiated Mellin
integrals converge at every \(s_0\).
The identity \(\mathcal M E_\zeta\psi=\zeta\,\mathcal M\psi\)
makes them vanish to order \(m\).
The window integrals are therefore the negatives
of the tail integrals just bounded.
The normalized jets (TB.61), the same-ordinate
Hilbert Gram blocks (TB.63), and the distinct-ordinate
vanishing Gram blocks now give the lower limit
\[
 \liminf_{\Lambda\to\infty}
       \operatorname{Tr}(D_{\Lambda,0}C_aC_a^*)
                 \geq\sum_\gamma m_\gamma|a(\gamma)|^2
 \tag{TP.32}
\]
for every compactly supported smooth multiplier \(a\).
To be explicit about the order of limits, first keep
any finite set of zero jets fixed, use their
asymptotic Gram matrix and their vanishing
\(Q_{\Lambda,0}\)-projections, and apply the finite
Bessel inequality.  Then increase the finite set.
Thus no closure of all jets in a fixed window is used.

\emph{3. The exact two-polar residual identity.}
Use the smooth raw \(P=P_\Lambda^p\), \(Q=Q_\Lambda^p\)
of (TU.1); they are distinct from the sharp
projections used in steps 1--2.
For a compact-frequency multiplier \(a\), use the
finite-pole inverse \(b=a/\zeta(1/2+i\tau)\),
the zero-jet projection \(P_{Z,L}\), and
\[
             W_L=S_L-P_{Z,L},\qquad
             T_L=U^{-1}C_bW_L .
 \tag{TP.33}
\]
These are precisely the constructions (TU.7)--(TU.14).
The pole of \(\zeta\) at \(s=1\) causes no pole of
\(1/\zeta\); the distribution poles in (TP.33)
are exactly its critical zeros.
The same polynomial cancellation proves the
Hilbert--Schmidt log-Sobolev bounds and
\(E_\zeta T_L=C_aW_L\).
The raw Fourier multiplier (TU.13) is the even one.
It gives the same reflected boundary prescription
and the same two cross estimates (TU.16).

Put \(A_L=PQPT_L\), and define its two normalized
constraint rows and its exactly constrained version by
\[
 r_+=I(A_L)/\sqrt{\Lambda},\qquad
 r_-=A_L(0)/\sqrt{\Lambda},\qquad
 A_L^0=A_L-b_{\Lambda,+}r_+-b_{\Lambda,-}r_- .
 \tag{TP.34}
\]
Both constraints of \(A_L^0\) vanish exactly.
Let \(S_+=1_{t\leq L}\), \(S_-=1_{t\geq-L}\).
The principal replacement for (TU.17) is the exact
identity
\[
 \begin{split}
 S_LE_\zeta T_L-E_\zeta A_L^0
 ={}&S_LE_\zeta(1-P)T_L
       +S_LE_\zeta(1-Q)PT_L\\
    &+S_+E_\zeta(1-P)QPT_L\\
    &+S_+E_\zeta b_{\Lambda,+}\,r_+
       +S_-E_\zeta b_{\Lambda,-}\,r_- .
 \end{split}
 \tag{TP.35}
\]
To verify every polar term, set \(g=QPT_L\).
The upper tail of \(E_\zeta Pg\) is
\(-I(Pg)e^{-t/2}\) and the lower tail of
\(E_\zeta g\) is \(-g(0)e^{t/2}/2\).
Expanding \(S_LE_\zeta T_L-E_\zeta Pg\)
as in (TU.17) consequently adds
\[
 r_+\,1_{t>L}e^{-(t-L)/2}
       +\tfrac12 r_-\,1_{t<-L}e^{(t+L)/2}
 \tag{TP.36}
\]
to its first three terms.
By (TP.5)--(TP.6), these two profile tails are
respectively
\(-1_{t>L}e^{-(t-L)/2}\) and
\(-\tfrac12 1_{t<-L}e^{(t+L)/2}\).
Adding the two corrections in (TP.34) therefore
gives exactly (TP.35).
The identities extend from smooth cores by the
graph estimates in (TP.3), just as in the
support-domain argument of AC-T39.

\emph{4. The two additional boundary rows have uniform bounds.}
We record this detail because a crude estimate by
\(\|T_L\|_2\) would introduce an unacceptable
power of \(L\) in the later frequency argument.
Up to \(O_M(c^{-M}\|T_L\|_2)\), the first row is
\[
             \Lambda^{-1/2}I(P^2T_L).
 \tag{TP.37}
\]
Indeed, in the difference
\(I(P(Q-1)PT_L)\), move \(Q-1\) to the normalized
smooth raw test \(p(u/\Lambda)/\sqrt{\Lambda}\).
Its Fourier tail outside the inner band has
every polynomial decay rate; Cauchy--Schwarz
proves the stated error.
After shifting the input endpoint \(L\) to zero,
the row (TP.37) is the integral against
\[
 \theta_2(t)=e^{t/2}p(e^t)^2
\]
of the kernel \(\kappa(t-s)\), for \(s\leq0\),
followed by the input contraction \(W_L\).
The Fourier transform of \(\theta_2\), with the
appropriate reflected sign, is analytic and
uniformly bounded on \(|\operatorname{Im}\tau|<1/4\):
at \(-\infty\) its defining integral decays at
least as \(e^{t/4}\), and its upper support is
fixed.  Thus this row is a reflected positive
half-line kernel of \(b\) with an additional
fixed analytic numerator.

For the second row, evaluating the band convolution
at zero gives, up to the same negligible error,
\[
          \frac2{\sqrt{\Lambda}}\,
                 I\bigl(p(\xi/\Lambda)\mathcal F T_L(\xi)\bigr).
 \tag{TP.38}
\]
For the error, move \(1-P\) to the inverse Fourier
transform of \(p(\xi/\Lambda)/\sqrt{\Lambda}\).
Its raw time tail has every polynomial rate.
Equation (TP.38) uses the identical test
\(\theta_1(t)=e^{t/2}p(e^t)\), and the Fourier-dual
inverse kernel from (TU.14).  It has the same
analytic-numerator description.

For a fixed compact frequency interval these two
rows converge in \(L^2((-\infty,0])\), with the
same harmless input truncations as in AC-T39.
For a frequency interval centered at height \(T\),
take the adaptive bump in AC-T40.
At \(|T|\geq10\) the zeta pole is outside its
disks; its proof applies literally to \(\zeta\).
The Euler center bound and the polynomial strip
bound hold there, and all zeros in the disks
are critical under \(\mathsf A\).
The analytic-numerator version of (FB.3) applied
to \(\theta_1,\theta_2\) proves a bound
\[
       \|r_+\|+\|r_-\|\leq C(2+|T|)^d
             +C_Mc^{-M}\|T_L\|_2 .
 \tag{TP.39}
\]
It also proves every fixed derivative or log-moment
bound on the uncontracted endpoint rows that is
needed here.  The input \(W_L\) is a contraction,
so no zero-jet Gram inverse or minimum zero gap
is introduced.  The two fixed output profiles
\(S_+E_\zeta b_{\Lambda,+}\) and
\(S_-E_\zeta b_{\Lambda,-}\), after endpoint
translation, have finite first log moment and
belong to \(H^s\) for every \(s<1/2\).
This follows from (TP.5): their uncut profiles
have exponential polar tails at one end and
rapid decay at the other, and the sharp half-line
cutoff has precisely the fractional endpoint
bound used in (TB.33).
Therefore both additional terms in (TP.35) obey
the same endpoint \(H^{1/4}\), \(L^2\), and first
log-moment column bounds as the original boundary
families.

\emph{5. Compact-frequency upper convergence.}
The first three terms of (TP.35) have the
graph-compact boundary descriptions proved in
(TU.15)--(TU.16a), with \(E_\zeta\) in place of
\(E_\chi\).  The second term uses (TP.3) to
exchange the lower endpoint for the upper one.
The preceding step gives the same description
for the two extra terms.
Each target boundary lies on the indicated
half-line.  AC-T43 approximates its fixed
Hilbert--Schmidt limit by principal \(S_0\)
profiles; AC-T42 replaces those profiles by
actual vectors in \(\mathcal R_\Lambda^0\).
For the smooth bulk \(A_L^0\), the last paragraph
of AC-T42 gives a graph approximation by that
same actual finite space.
Indeed the uncorrected bulk has the arbitrary
gap estimate of AC-T38, and the two correction
profiles in (TP.34) have the gap estimates
(TP.7); their finite constraint errors are
removed by the exact map (TP.9).
For this fixed frequency interval all their
column norms grow at most polynomially in \(L\),
so the arbitrary powers of \(c^{-1}\) suffice.
The two input-outside-window families
\(S_LC_a(1-S_L)\) are the same Schwartz quadrant
kernels as in (TU.21), and AC-T43 absorbs
them as well.
It follows that
\[
       \|(1-Q_{\Lambda,0})S_LC_a(1-P_{Z,L})\|_2
                                          \longrightarrow0 .
 \tag{TP.40}
\]
The exact positive distance formula is
\(\operatorname{Tr}(D_{\Lambda,0}C_aC_a^*)
 =\|(1-Q_{\Lambda,0})S_LC_a\|_2^2\).
On the finite zero-jet input space its upper
limit is the sum in (TP.32), by (TU.8).
Equation (TP.40) handles the orthogonal input
complement.  This proves compact-frequency
equality, and polarization gives vague
convergence of the positive measures
\[
 d\mu_{\Lambda,0}(\tau)
   =\frac{2L-\sum_j|\widehat{S_Lz_j}_+(\tau)|^2}{2\pi}\,d\tau,
 \qquad z_j=E_0(G_\Lambda^{\rm pr})^{-1/2}e_j .
 \tag{TP.41}
\]
Their positivity is Bessel's inequality, exactly
as in AC-T30.

\emph{6. Uniform frequency control and the original tests.}
We check that the principal modifications preserve
the quantitative threshold, rather than relying
on compact-frequency convergence alone.
Use AC-T40's bump equal to one on \([T,T+1]\),
with the same fixed derivative and moment budgets
as in AC-T41.  The reciprocal \(1/\zeta\) has
only the finitely many critical zero poles in
that window.  The local zero count is
\(N\leq C\log(2+|T|)\), and the finite-pole
construction gives
\[
       \|T_L\|_{2,H^j}
          \leq C(2+|T|)^d(1+L)^{N+d}
 \tag{TP.42}
\]
for every fixed derivative budget.
All uncorrected boundary families have
height-polynomial target norms by AC-T40 and
the \(E_\zeta\) multiplier and derivative bounds.
The two new rows have the same bounds by
(TP.39).  Right composition by the exact input
contraction retains these bounds.
AC-T43 therefore absorbs every boundary within
any prescribed fixed tolerance once
\(L\geq C(2+|T|)^D\), with a fixed exponent \(D\).
At that same threshold, the negligible errors
in (TP.39), the cross estimates, and all
constrained-bulk graph errors tend uniformly
to zero: their logarithms are bounded by
\[
       C\log(2+|T|)+(N+C)\log(1+L)-c_0L ,
 \tag{TP.43}
\]
after choosing a fixed sufficiently strong gap
power.  Since \(N=O(\log(2+|T|))\), the last
negative term dominates for the stated polynomial
threshold in \(L\).
Thus the orthogonal-complement term in (TP.40)
has squared norm at most one above that threshold.
The zero-jet term is bounded by its rank \(N\),
because the bump has norm at most one.
It follows that
\[
 \mu_{\Lambda,0}([T,T+1])
       \leq C(2+|T|)^A
              \quad\hbox{for all sufficiently large }\Lambda,
 \tag{TP.44}
\]
with fixed \(A,C\).
Below the threshold use the elementary density
bound \(\mu_{\Lambda,0}([T,T+1])\leq L/\pi\).
For bounded heights use the already proved local
upper convergence and then that same elementary
bound on the remaining bounded range of \(\Lambda\).
This proves (TP.44) in all cases.

The polynomial unit-interval bound makes the
tails against any Schwartz function uniformly
summable.  Compact-frequency convergence thus
extends to all Schwartz multipliers by truncation.
In particular it applies to \(\widehat h_+\)
for every original \(h\in C_c^\infty(\mathbb R_+)\).
The finite-rank Fourier trace identity for
(TP.41), followed by (TP.29), gives (TP.19).
All cutoffs throughout this proof are the original
full even prolate cutoff with the two explicitly
stated principal constraints.
\end{proof}

